\section{The wing constant, and existence}\label{sec:exist}

The log discrete local variance $\lambda=(\lambda^2,\dots,\lambda^M)$,
$\lambda^j=\log\dlv^j$, ranges over all of $\R^{N(M-1)}$, so positivity is
automatic; the front density has its own chart \eqref{eq:softmax}. By Theorem~\ref{thm:pointwise} the model
is the stochastic differential equation
\begin{equation}\label{eq:abstract}
  \dd\lambda_t=\Bigl[-\ip{\eta_t}{b_t}-\tfrac12\abs{b_t}^{2}\Bigr]\dd t
   +b_t\,\dd W_t ,
\end{equation}
whose only non-elementary coefficient is $\eta=\psi/q$ from \eqref{eq:volsolve}
and \eqref{eq:eta}. Its denominator is small in the wings, and the question is
whether $\eta$ can blow up there.

The argument has the structure of Wissel's proof of his Theorem~3.5
\cite{Wissel2007}: log coordinates, induction over expiries, and local existence
followed by a bound that rules out explosion. Its one non-standard ingredient is an a priori bound, independent of the state, on the
relative volatility of the calendar spread at the second expiry. Wissel could not
obtain this bound in general and imposes his condition (3.24) instead (his
Remark~1 following the proof, p.~18); in our coordinates the quantity is $\eta+b$
(Appendix~\ref{app:wissel}). We obtain it from an exact identity for the
volatility recursion and a known inequality for inverses of M-matrices.

\subsection{The volatility recursion}\label{sec:walk}\label{sec:master}

Write $P:=A_j(\dlv^j)^{-1}$, so that $q^j=Pq^{j-1}$.

\begin{corollary}[Positivity, unit column sums, absorbing ends]\label{cor:pos}
$P\ge0$ entrywise, $\one^{\!\top}P=\one^{\!\top}$, so $P_{im}\in[0,1]$, and the outer
columns of $P$ are $e_0$ and $e_{N+1}$.
\end{corollary}

\begin{proof}
By \eqref{eq:operator}, $A_j$ has non-positive off-diagonal entries, and in each
column the diagonal entry exceeds the sum of the moduli of the off-diagonal ones by
exactly one; so $A_j$ is a non-singular M-matrix and $P\ge0$
\cite{BermanPlemmons1994}. The column sums follow from $\one^{\!\top}A_j=\one^{\!\top}$
(Lemma~\ref{lem:null}), and the outer columns of $A_j$ are $e_0$ and $e_{N+1}$. The
first two statements are also Theorem~3.3 of \cite{BuehlerRyskin2016}, attributed
there to \cite{AndreasenHuge2011}.
\end{proof}

That is all that is needed about $P$ itself. What we want next is a recursion for
$\eta$ alone: one step per expiry, each adding a controlled amount, so that the
bound telescopes along the ladder. Dividing \eqref{eq:volsolve} by $q^j$ gives
\begin{equation}\label{eq:divided}
  \eta^{j,i}\;=\;\frac{(P\psi^{j-1})_i}{q^{j,i}}
   \;-\;\frac{\bigl(P\,G_j(q^j)\beta^j\bigr)_i}{q^{j,i}} ,
\end{equation}
the first term carrying $\eta^{j-1}$ and the second carrying $b^j$.

The first is an average. Writing $\psi^{j-1}=q^{j-1}\!\circ\eta^{j-1}$, its weights
$P_{ik}q^{j-1,k}/q^{j,i}$ are non-negative because $P\ge0$ and sum to one because
$q^j=Pq^{j-1}$, so it contributes at most $\lVert\eta^{j-1}\rVert_\infty$; they are
the matrix $\mathcal A$ of \eqref{eq:mastercomp}.

The source is the difficulty. Since $G_j$ is a second difference, row $i$ carries
the neighbouring densities $q^{j,i\pm1}$ as well as $q^{j,i}$, so dividing by $q^{j,i}$
leaves the ratios
\begin{equation}\label{eq:nbr}
  \frac{q^{j,i-1}}{q^{j,i}}
  \qquad\text{and}\qquad
  \frac{q^{j,i+1}}{q^{j,i}} ,
\end{equation}
which are unbounded in the wings, where the density decays fastest; nothing in
\eqref{eq:volsolve} bounds them.

They are, however, an artefact of writing the row out. Since
$\one^{\!\top}G_j(q)=0$ by Lemma~\ref{lem:null}, row $i$ of $G_j(q)\dlv^j$ is the
net outflow of node $i$, and by \eqref{eq:affine} that net is
\begin{equation}\label{eq:netknown}
  G_j(q^j)\,\dlv^j\;=\;A_j(\dlv^j)q^j-q^j\;=\;q^{j-1}-q^{j},
\end{equation}
which carries no ratios at all: the individual fluxes are wing quantities, their
net is not. For $b^j\equiv c$ this settles the matter, since $\beta^j=c\,\dlv^j$
and $Pq^{j-1}=q^j$ give
\[
  \frac{\bigl(P\,G_j(q^j)\beta^j\bigr)_i}{q^{j,i}}
  =c\,\frac{\bigl(q^{j}-Pq^{j}\bigr)_i}{q^{j,i}}
  =c\Bigl(1-\frac{(Pq^{j})_i}{q^{j,i}}\Bigr),
\]
the two ratios of \eqref{eq:nbr} having collapsed into the single propagation
ratio $(Pq^j)_i/q^{j,i}$. A general $b$ only modulates the density. Let
$E\in\R^{N\times(N+2)}$ restrict to the interior nodes, so $EE^{\!\top}=I_N$, and
note that $A_j=I+\Gh_j\,\mathrm{diag}(\dlv^j)E$. Writing
$E^{\!\top}(q\circ b)$ for the vector equal to $q^ib^i$ at the interior nodes and
zero at the outer ones, \eqref{eq:factor} and the commuting of diagonal matrices
give
\begin{equation}\label{eq:modulated}
  G_j(q)\beta^j
  =\Gh_j\,\mathrm{diag}(\dlv^j)E\,E^{\!\top}(q\circ b^j)
  =\bigl(A_j(\dlv^j)-I\bigr)E^{\!\top}(q\circ b^j),
\end{equation}
again the operator's own increment, applied to $E^{\!\top}(q\circ b^j)$ in place of
$q$. The proposition is what this amounts to.

\begin{proposition}\label{prop:master}
With $b^j=\beta^j/\dlv^j$ and $\circ$ the entrywise product,
\begin{equation}\label{eq:master}
  \psi^{j}\;=\;P\,\psi^{j-1}\;-\;(I-P)\,E^{\!\top}\bigl(q^{j}\!\circ b^{j}\bigr),
\end{equation}
equivalently, with the entrywise non-negative
\[
  \mathcal A:=\mathrm{diag}(q^{j})^{-1}P\,\mathrm{diag}(q^{j-1}),
  \qquad
  \mathcal B:=\mathrm{diag}(q^{j})^{-1}P\,E^{\!\top}\mathrm{diag}(Eq^{j}),
\]
\begin{equation}\label{eq:mastercomp}
  \eta^{j}\;=\;\mathcal A\,\eta^{j-1}\;+\;\bigl(\mathcal B-E^{\!\top}\bigr)b^{j} ,
\end{equation}
where $\mathcal A\one=\one$, so the transport is an averaging operator, and the
$i$-th row of $\mathcal B$ sums to
$(Pq^{j})_i/q^{j,i}-\one_{\{i\in\{0,N+1\}\}}$.
\end{proposition}

\begin{proof}
By \eqref{eq:volsolve}, $\psi^j=P\psi^{j-1}-P\,G_j(q^j)\beta^j$. Applying $P$ to
\eqref{eq:modulated} at $q=q^j$ and using $PA_j=I$ gives
$P\,G_j(q^j)\beta^j=(I-P)E^{\!\top}(q^j\!\circ b^j)$, which is \eqref{eq:master}.
Dividing it by $q^{j,i}$ gives \eqref{eq:mastercomp}: the rows of $\mathcal A$ sum to
$(Pq^{j-1})_i/q^{j,i}=1$, and for $\mathcal B$ the terms $k\in\{0,N+1\}$ missing from
$(Pq^j)_i$ vanish for interior $i$ and equal $q^{j,i}$ at $k=i$ for outer $i$, by
Corollary~\ref{cor:pos}.
\end{proof}

So $\eta^j$ is an average of $\eta^{j-1}$ plus a source in $b^j$, and the march
across expiries is additive.

\subsection{The wing constant}\label{sec:bound}

That single surviving ratio, $(Pq^j)_i/q^{j,i}$, is what has to be bounded. Since
$q^j=Pq^{j-1}$ it is
\begin{equation}\label{eq:P2ratio}
  \frac{(Pq^{j})_i}{q^{j,i}}=\frac{(P^{2}q^{j-1})_i}{(Pq^{j-1})_i} ,
\end{equation}
so what is needed is how much larger than $P$ the square $P^{2}$ can be entrywise;
and since both sides are linear and non-negative in $q^{j-1}$, the case
$q^{j-1}=e_m$ shows that an entrywise bound is not merely sufficient but
necessary. Write
\begin{equation}\label{eq:rows}
  R^{j,i}:=\one_{\{1\le i\le N\}}+\frac{(Pq^j)_i}{q^{j,i}}-\one_{\{i\in\{0,N+1\}\}},
  \qquad
  \Theta_j:=\max_i R^{j,i},
  \qquad
  \Xi:=\sum_{r=1}^{N}P_{rr}\le N ,
\end{equation}
and $H_j:=\max_i\abs{\eta^{j,i}}$. Here $R^{j,i}$ is the row bound of
\eqref{eq:mastercomp} at node $i$, and $\Theta_j$, the \emph{wing constant}, is its
size across the slab. The theorem shows $\Theta_j$ is bounded by the grid alone,
uniformly in the state; the bound rests on a single inequality for inverses of
M-matrices.

\begin{lemma}[Squared propagator]\label{lem:pathproduct}\label{cor:p2p}
For all nodes $i,r,m$,
\begin{equation}\label{eq:summand}
  P_{ir}\,P_{rm}\;\le\;P_{im}\,P_{rr},
\end{equation}
and consequently, for all $i,m$,
\begin{equation}\label{eq:p2p}
  \bigl(P^{2}\bigr)_{im}\le
  \begin{cases}
    \Xi\,P_{im}, & 1\le i\le N,\\
    (1+\Xi)\,P_{im}, & i\in\{0,N+1\}.
  \end{cases}
\end{equation}
\end{lemma}

\begin{proof}
$P$ is the inverse of an M-matrix by Corollary~\ref{cor:pos}, and every inverse
M-matrix is a path product matrix, $a_{ij}a_{jk}\le a_{ik}a_{jj}$ for all
$i,j,k$ \cite[Theorem~4.5]{JohnsonSmith2011}; \eqref{eq:summand} is the case
$(i,j,k)=(i,r,m)$. Summing it over $r$ gives \eqref{eq:p2p}: an outer $r$
contributes only if $r=i$, since column $r$ of $P$ is $e_r$, and then contributes
$P_{im}$.
\end{proof}

\begin{theorem}\label{thm:wing}
For every $\dlv\in(0,\infty)^N$,
\begin{equation}\label{eq:step}
  H_j\;\le\;H_{j-1}+\Theta_j\lVert b^j\rVert_\infty ,
\end{equation}
and
\begin{equation}\label{eq:thetabound}
  \Theta_j\le1+\Xi\le N+1,
  \qquad
  \Theta^{*}:=\sup_{\dlv}\sum_{j=2}^{M}\Theta_j\le(M-1)(N+1).
\end{equation}
\end{theorem}

\begin{proof}
Since $q^j=Pq^{j-1}\ge0$, Lemma~\ref{lem:pathproduct} gives
$(Pq^j)_i=\sum_m(P^2)_{im}q^{j-1,m}\le\Xi\,q^{j,i}$ for interior $i$ and
$\le(1+\Xi)q^{j,i}$ for outer $i$, so $R^{j,i}\le1+\Xi$ in both cases. Finally,
$\mathcal A\ge0$ with $\mathcal A\one=\one$ gives
$\lVert\mathcal A\eta^{j-1}\rVert_\infty\le H_{j-1}$, while row $i$ of
$(\mathcal B-E^{\!\top})b^{j}$ has modulus at most $\lVert b^j\rVert_\infty R^{j,i}$
by the row sums of Proposition~\ref{prop:master}.
\end{proof}

\subsection{Existence}\label{sec:existence}

Iterating \eqref{eq:step} from the front slice,
\begin{equation}\label{eq:apriori}
  \sup_{j,i}\abs{\eta^{j,i}_{t}}
  \;\le\;\lVert\eta^1_t\rVert_\infty+(M-1)(N+1)\,\lVert b\rVert_\infty ,
\end{equation}
and iterating \eqref{eq:master} shows that $\eta^j$ is affine in $\eta^1$, with
coefficients depending on $\lambda$ only through $P_2,\dots,P_j$ and $q^2,\dots,q^j$.
The only $t$ on the right of \eqref{eq:apriori} is inside $\eta^1_t$, which enters
linearly, so it has to be integrable in time rather than bounded.

\begin{assumption}[Integrable front volatility]\label{ass:F}
For a horizon $T\le T_1$,
\begin{equation}\label{eq:F}
  \int_0^{T}\lVert\eta^{1}_t\rVert_\infty\,\dd t\;<\;\infty\qquad\text{almost surely.}
\end{equation}
\end{assumption}

\begin{theorem}[Existence]\label{thm:exist}
Let $T<T_1$, and let the front volatility $\nu=\nu(t,z)$ in \eqref{eq:softmax} and
the volatility $b=b(t,z,\lambda)$ be bounded, $\abs{\nu}\le B_\nu$,
$\abs{b}\le B$, and locally Lipschitz in $(t,z,\lambda)$. Then the front chart
\eqref{eq:softmax}--\eqref{eq:frontdrift} together with \eqref{eq:abstract} has a
unique strong solution $(z,\lambda)$ on $[0,T]$, and
\begin{equation}\label{eq:ellbounds}
  0\;<\;\inf_{t\le T}\min_{j,i}\dlv^{j,i}_{t}
  \;\le\;\sup_{t\le T}\max_{j,i}\dlv^{j,i}_{t}\;<\;\infty
  \qquad\text{almost surely.}
\end{equation}
No condition is imposed on $\dlv$, on the wings, or on the number of factors.
\end{theorem}

\begin{proof}
The coefficients are locally Lipschitz: $q^1=\mathrm{softmax}(z)$ and
$\dlv=e^\lambda$ are positive, so each $A_j$ is a non-singular M-matrix and
everything except $\nu$ and $b$ is smooth in $(z,\lambda)$. Adding time as a state
variable to cover the time dependence, \cite[Theorem~V.38, p.~303]{Protter2005}
gives a unique solution up to an explosion time $\zeta$, with
$\limsup_{t\to\zeta}\lVert(z_t,\lambda_t)\rVert=\infty$ on $\{\zeta<\infty\}$.
The drift of $\lambda$ is bounded by \eqref{eq:apriori}, since
$\lVert\eta^1\rVert_\infty=\lVert\mathrm{dev}\rVert_\infty\le2B_\nu$, and the drift
$m$ of $z$ is bounded by $2B_\nu^2$; let $D$ bound both. On $[0,\zeta\wedge T)$
the drift part of $(z,\lambda)$ is therefore at most $DT$, and the noise part,
having integrand bounded by $B_\nu+B$, is a martingale bounded in $L^2$, hence
convergent as $t\to\zeta\wedge T$. So $(z,\lambda)$ stays bounded on
$[0,\zeta\wedge T)$, which rules out $\zeta\le T$. The solution exists on $[0,T]$,
$\lambda$ is bounded there, and \eqref{eq:ellbounds} follows from $\dlv=e^\lambda$.
\end{proof}

\begin{remark}[The front expiry is included]\label{rem:horizon}
At $T_1$ the front settles on a node (Proposition~\ref{prop:collapse}), so $z$
itself runs off and Theorem~\ref{thm:exist} does not apply on $[0,T_1]$ directly.
It applies on every $[0,T']$ with $T'<T_1$, and the $\lambda$-part of the proof
only needs the drift bound to be integrable in time, which is
Assumption~\ref{ass:F} with $T=T_1$. Under it $\lambda$ stays bounded on $[0,T_1)$
and has a limit at $T_1$, so \eqref{eq:ellbounds} holds with $T=T_1$. The settling
front has $\lVert\eta^1_t\rVert_\infty\le C(T_1-t)^{-p}$, integrable for $p<1$, so
the code-book reaches its first expiry with every gridded contract settling at its
payoff.
\end{remark}

\subsection{Beyond the first expiry}\label{sec:patching}

At $T_1$ the front expires and the next expiry becomes the front. The settling of
Proposition~\ref{prop:collapse} makes this switch clean: the spot sits on a node, and the
new front density is a column of the propagator. Indeed, if $q^l_{T_l}=e_k$ then the
code-book $A_{l+1}(\dlv^{l+1})\,q^{l+1}=q^l$ gives
\begin{equation}\label{eq:handover}
  q^{l+1}_{T_l}=P_{l+1}\,e_k ,\qquad P_{l+1}:=A_{l+1}\bigl(\dlv^{l+1}_{T_l}\bigr)^{-1},
\end{equation}
the spread of a unit mass at the settled spot by one step of the code-book. To restart the
front chart we need this density to be strictly positive.

\begin{lemma}[The new front]\label{lem:newfront}
Let $P=A_j(\dlv^j)^{-1}$ with $\dlv^j\in(0,\infty)^N$. If $k$ is an interior node, then
$(Pe_k)^i>0$ for every node $i$ and $K^{\!\top}Pe_k=K^k$. If $k\in\{0,N+1\}$, then
$Pe_k=e_k$.
\end{lemma}

\begin{proof}
Order the nodes with the interior first. The outer columns of $A_j$ are $e_0$ and
$e_{N+1}$ (Corollary~\ref{cor:pos}), so
\[
  A_j=\begin{pmatrix}A_{II}&0\\A_{BI}&I\end{pmatrix},\qquad
  P=\begin{pmatrix}A_{II}^{-1}&0\\-A_{BI}A_{II}^{-1}&I\end{pmatrix}.
\]
By \eqref{eq:operator}, $A_{II}$ is tridiagonal with strictly negative off-diagonal
entries, hence an irreducible non-singular M-matrix, so $A_{II}^{-1}>0$ entrywise
\cite{BermanPlemmons1994}. The block $-A_{BI}\ge0$ has a positive entry in the row of
$K^0$ (column $K^1$) and in the row of $K^{N+1}$ (column $K^N$), so the outer entries of an
interior column of $P$ are positive as well. The identity $K^{\!\top}P=K^{\!\top}$ follows
from $K^{\!\top}A_j=K^{\!\top}$ (Lemma~\ref{lem:null}), and the outer columns are
Corollary~\ref{cor:pos}.
\end{proof}

\begin{corollary}[Existence on the whole ladder]\label{cor:patch}
Let $b$ be bounded and locally Lipschitz. On each interval $[T_{l-1},T_l]$,
$l=1,\dots,M$, let the current front density $q^l$ follow the chart \eqref{eq:softmax}
with the settling volatility \eqref{eq:nufront}, where $\tau=T_l-t$ and
$p\in[\tfrac12,1)$. Then the model has a solution on $[0,T_M]$, unique given these
choices, such that
\begin{enumerate}[label=(\roman*),leftmargin=2em,itemsep=1pt]
\item every gridded price is a continuous martingale up to its expiry;
\item at every $T_l$ the spot sits on a node $K^{k_l}$ and every contract expiring at
  $T_l$ settles at its payoff;
\item if $k_l$ is interior, the model continues on $[T_l,T_{l+1}]$ from the finite state
  $z=\log q^{l+1}_{T_l}$ and $\lambda^{l+2},\dots,\lambda^M$ at $T_l$; if $k_l$ is an
  outer node, the spot is absorbed there: $q^j_t=e_{k_l}$ for all $j>l$ and $t\ge T_l$.
\end{enumerate}
\end{corollary}

\begin{proof}
By induction on $l$. On $[0,T_1]$ this is Proposition~\ref{prop:collapse} with
Theorem~\ref{thm:exist} and Remark~\ref{rem:horizon}, which also give the limit
$\lambda_{T_1}$, finite. Suppose the solution exists on $[0,T_l]$ with $\lambda_{T_l}$
finite and $q^l_{T_l}=e_k$.

If $k$ is interior, \eqref{eq:handover} and Lemma~\ref{lem:newfront} give a strictly
positive new front, so $z=\log q^{l+1}_{T_l}$ is finite, and
$K^{\!\top}q^{l+1}_{T_l}=K^k=S_{T_l}$, so the spot is continuous. On $[T_l,T_{l+1}]$ the
ladder is that of Sections~\ref{sec:drift}--\ref{sec:exist} with front $q^{l+1}$ and
expiries $T_{l+1}<\dots<T_M$, started from the $\mathcal F_{T_l}$-measurable state
$(z,\lambda^{l+2},\dots,\lambda^M)$. Theorem~\ref{thm:exist} applies with this
$\mathcal F_{T_l}$-measurable initial condition, the solution of
\cite[Theorem~V.38]{Protter2005} depending continuously on its starting point,
and with Proposition~\ref{prop:collapse} and Remark~\ref{rem:horizon} gives a unique
strong solution on $[T_l,T_{l+1}]$ that settles at $T_{l+1}$ with $\lambda_{T_{l+1}}$
finite. The densities $q^j$, $j>l+1$, are given by the same recursion on both sides of
$T_l$ with continuous inputs, so they are continuous at $T_l$.

If $k$ is an outer node, Lemma~\ref{lem:newfront} gives $q^j_{T_l}=e_k$ for every $j>l$.
The constant process $q^j_t\equiv e_k$ satisfies every recursion, since $A_je_k=e_k$, and
is a martingale, so all prices are constant and settle at their payoffs; we freeze
$\lambda$ at $\lambda_{T_l}$, on which no price then depends.

Uniqueness on each interval gives uniqueness on $[0,T_M]$.
\end{proof}
